\def\urlSite{https://brevet.initiation.aero}

\ifdefined\activerhandout 
\documentclass[12pt,aspectratio=1610,handout]{beamer}
\else
\documentclass[12pt,aspectratio=1610]{beamer}
\fi

\usepackage[
    backend=biber,
    %style=authoryear-icomp,
    style=numeric,
    %style=ieee,
    sortlocale=fr_FR,
    natbib=true,
    url=true, 
    doi=true,
    eprint=true
]{biblatex}
	\addbibresource{commun/biblio.bib}
	\addbibresource{01-EtudeAeronefs/biblio.bib}
	\addbibresource{02-Navigation/biblio.bib}
	\addbibresource{03-Meteo/biblio.bib}
	\addbibresource{04-Aerodynamique/biblio.bib}
	\addbibresource{05-Histoire/biblio.bib}
	\addbibresource{Annexes/biblio.bib}
	\addbibresource{Annexes/friseChronologique/biblio.bib}
	\addbibresource{Annexes/friseChronologique/extraitsVideo/biblio.bib}

%\usepackage[french]{babel} %=> erreur avec les tikz du moteur
\usepackage[T1]{fontenc}
\usepackage[utf8]{inputenc}
\usepackage{lmodern}
\usepackage{hyperref}
\usepackage{smartdiagram}
\usepackage{xcolor}
\usepackage{tikz}
\usepackage{animate}
\usepackage{tikzpeople}
\usepackage{appendixnumberbeamer}
\usepackage[labelformat=empty]{caption}
%\usepackage{pictochrono}
\usepackage{fontawesome5}
\usepackage{awesomebox}
\usepackage{siunitx}
\usepackage{icomma}
\usepackage{physics}
\usepackage[pdftex,outline]{contour} % glow around text
\usepackage{ifthen}
\usepackage{fp}
\usepackage{textcomp}
\usepackage{gensymb}

\usetheme{Warsaw}
\setbeamertemplate{page number in head/foot}[totalframenumber]

\newcommand{\anglais}[1]{(\textit{\color{blue}#1})}
\newcommand{\legende}[2]{\caption{#1\ifdefined\activerbibliobeamer \cite{#2}\fi}}
\newcommand{\histoire}[1]{\begin{awesomeblock}{2pt}{\faBook}{black!75}#1\end{awesomeblock}}
\newcommand{\info}[1]{\begin{awesomeblock}{2pt}{\faInfoCircle}{blue!75}#1\end{awesomeblock}}
\newcommand{\question}[1]{\begin{awesomeblock}{2pt}{\faQuestionCircle}{black!75}#1\end{awesomeblock}}
\newcommand{\alerte}[1]{\begin{awesomeblock}{2pt}{\faExclamationCircle}{red!75}#1\end{awesomeblock}}
\newcommand{\astuce}[1]{\begin{awesomeblock}{2pt}{\faLightbulb}{black!75}#1\end{awesomeblock}}
\newcommand{\exemple}[1]{\begin{awesomeblock}{2pt}{\faSearch}{black!75}#1\end{awesomeblock}}
\newcommand{\definitionAConnaitre}[1]{\begin{awesomeblock}{2pt}{\faCog}{black!75}#1\end{awesomeblock}}

%\newcommand{\qmcBia}[8]{
%%1 : titre slide
%%2 : numéro de la bonne réponse
%%3 : Inititulé de la question
%%4, 5, 6, et 7 : propositions de réponse
%%8 : explication
%\begin{frame}{#1}
%\begin{awesomeblock}{2pt}{\faQuestion}{black!75}
%#3
%	\begin{enumerate}
%	\ifnum#2=1
%		\only<1>{\item #4}
%		\only<2>{\item \textbf{#4}}
%	\else
%		\item #4
%	\fi
%	\ifnum#2=2
%		\only<1>{\item #5}
%		\only<2>{\item \textbf{#5}}
%	\else
%		\item #5
%	\fi
%	\ifnum#2=3
%		\only<1>{\item #6}
%		\only<2>{\item \textbf{#6}}
%	\else
%		\item #6
%	\fi
%	\ifnum#2=4
%		\only<1>{\item #7}
%		\only<2>{\item \textbf{#7}}
%	\else
%		\item #7
%	\fi
%	\end{enumerate}
%	\pause
%\end{awesomeblock}
%#8
%\end{frame}
%}

\newcommand{\qmcBia}[8]{
%1 : titre slide
%2 : numéro de la bonne réponse
%3 : Inititulé de la question
%4, 5, 6, et 7 : propositions de réponse
%8 : explication
\begin{frame}{#1}
\begin{awesomeblock}{2pt}{\faQuestion}{black!75}
#3
	\begin{enumerate}
	\item #4
	\item #5
	\item #6
	\item #7
	\end{enumerate}
\end{awesomeblock}
\end{frame}

\begin{frame}{#1}
\begin{awesomeblock}{2pt}{\faQuestion}{black!75}
#3
	\begin{enumerate}
	\ifnum#2=1
		\item \textbf{#4}
	\else
		\item #4
	\fi
	\ifnum#2=2
		\item \textbf{#5}
	\else
		\item #5
	\fi
	\ifnum#2=3
		\item \textbf{#6}
	\else
		\item #6
	\fi
	\ifnum#2=4
		\item \textbf{#7}
	\else
		\item #7
	\fi
	\end{enumerate}
\end{awesomeblock}
#8
\end{frame}
}

\subtitle{BIA - Brevet d'Initiation Aéronautique}
%\author[BIA]{Clément \textsc{Vermot-Desroches}}
\author[BIA]{\urlSite}
\institute{Collège Aliénor d'Aquitaine\\Martignas-sur-Jalle}
%\date{\today}
\date{Année 2026-2027}

\AtBeginSection[]
{
    \begin{frame}
        %\frametitle{Table of Contents}
        \tableofcontents[currentsection]
    \end{frame}
}

\AtBeginSubsection[]
{
    \begin{frame}
        %\frametitle{Table of Contents}
        \tableofcontents[currentsection,currentsubsection]
    \end{frame}
}


\title[Séance 0]{Séance 0 \\ Réunion de présentation}

\begin{document}
 \begin{frame}
 \titlepage
 \end{frame}
 
 \begin{frame}
 \tableofcontents
 \end{frame}
 
\section{Le BIA}
\begin{frame}{Le BIA}

Le Brevet d'initiation aéronautique (BIA) :
\begin{itemize}
\item diplôme français correspondant attestant d'un culture générale du domaine de l'aéronautique
\item délivré conjointement par le ministère de la Transition écologique et solidaire, chargé des Transports, et par celui de l'Éducation nationale ainsi que par la Confédération nationale des Fédérations aéronautiques et sportives
\end{itemize}
\end{frame}

\begin{frame}{Le BIA}

Enseignement articulé autour de 5 thèmes :
\begin{enumerate}
\item Météorologie et aérologie
\item Aérodynamique, aérostatique et principes du vol
\item Étude des aéronefs et des engins spatiaux
\item Navigation, réglementation, sécurité des vols
\item Histoire et culture de l'aéronautique et du spatial
\end{enumerate}
\end{frame}

\begin{frame}{Format proposé au collège}
\begin{itemize}
	\item Séances de l'atelier BIA le jeudi (12h45-14h)
	\item 1 à 2 sorties proposées sur l'année hors temps scolaire (visite du CAEA, aéroclub...)
\end{itemize}

Compte tenu du nombre limité de séances, l'implication personnelle de chaque élève sera déterminante.

\end{frame}

\section{L'examen}
\begin{frame}{L'examen}
\begin{itemize}
	\item Examen national
	\item QCM
	\item Épreuve obligatoire de 100 questions (20 par thème), 2h30
	\item Une épreuve facultative d'anglais (20 questions), 30 minutes
	\item Examen organisé sur la secondes quinzaine du mois du mai
	\item Le collège devrait être centre d'examen, comme l'an dernier
	\item Mentions :
	\begin{itemize}
		\item 10/20 : admis
		\item 12/20 : assez bien
		\item 14/20 : bien
		\item 16/20 : très bien
	\end{itemize}
	\item en 2024, environ 21000 candidats, 81\% d'admis
\end{itemize}
\end{frame}

\section{Et après ?}
\begin{frame}
Le BIA constitue une première étape dans le monde de l'aéronautique. 

Il atteste d'un socle minimum de connaissance dans le domaine aéronautique et aérospatial. 

La plupart des fédérations aéronautiques allouent des bourses majorées aux jeunes titulaires d'un BIA qui souhaiteraient préparer une licence de pilotage (avion, ULM, planeur...).

Le BIA est également un diplôme différenciant sur le marché du travail, pour les entreprises avec une composante aéronautique ou spatiale.
\end{frame}

\end{document}

